% ============================================================
%  Basado en la plantilla de la Práctica 3 (mismos colores/estilos)
% ============================================================
\documentclass[11pt,a4paper]{article}

% ---------- Idioma y codificación ----------
\usepackage[utf8]{inputenc}
\usepackage[T1]{fontenc}
\usepackage[spanish,es-noshorthands]{babel}

% ---------- Tipografía y layout ----------
\usepackage{lmodern}
\usepackage[sfdefault,light]{roboto}
\usepackage{geometry}
\geometry{a4paper, top=3.1cm, bottom=2.8cm, left=2.4cm, right=2.4cm, headheight=1.6cm}
\usepackage{ragged2e}
\usepackage{enumitem}
\usepackage{graphicx}
\usepackage{float}
\usepackage{booktabs}
\usepackage{array}
\usepackage[most]{tcolorbox}
\usepackage{tikz}
\usetikzlibrary{shadows,calc}
\usepackage{titlesec}
\usepackage{fancyhdr}
\usepackage{fontawesome5}
\usepackage[colorlinks=true,breaklinks=true]{hyperref}
\usepackage{caption}
\usepackage{microtype}
\usepackage{xurl}

% ---------- Paleta de colores (tema "coolnight" de tu terminal Kitty) ----------
\definecolor{cNavy}{HTML}{011423}
\definecolor{cTealDeep}{HTML}{064B5A}
\definecolor{cTeal}{HTML}{097289}
\definecolor{cBlue}{HTML}{0FC5ED}
\definecolor{cCyan}{HTML}{0FC5ED}
\definecolor{cMint}{HTML}{44FFB1}
\definecolor{cPurple}{HTML}{A277FF}
\definecolor{cRed}{HTML}{E52E2E}
\definecolor{cBg}{HTML}{EAF6F8}
\definecolor{cText}{HTML}{01111E}

\hypersetup{
  linkcolor=cTealDeep,
  citecolor=cTealDeep,
  urlcolor=cBlue,
  pdftitle={Práctica 5 -- Reglas de Asociación con Apriori},
  pdfauthor={Miguel Deza}
}

\color{cText}
\emergencystretch=2.5em

% ---------- Encabezado / pie de página ----------
\pagestyle{fancy}
\fancyhf{}
\renewcommand{\headrulewidth}{0pt}
\renewcommand{\footrulewidth}{0pt}
\fancyhead[L]{\small\color{cTealDeep}\textbf{\faGraduationCap\ Maestría en Ciencia de Datos}}
\fancyhead[R]{\small\color{cTeal}Práctica 5 -- Reglas de Asociación}
\fancyfoot[C]{%
  \begin{tikzpicture}[remember picture]
    \node[circle, fill=cTealDeep, minimum size=1.7em, inner sep=0pt] (pg) {\color{white}\small\thepage};
  \end{tikzpicture}%
}
\renewcommand{\headrule}{{\color{cCyan}\hrule height 1.4pt}\vspace{1pt}{\color{cBlue}\hrule height 0.5pt}}

% ---------- Estilo de secciones ----------
\titleformat{\section}
  {\normalfont\Large\bfseries\color{cTealDeep}}
  {\tikz[baseline=-0.6ex]{\node[fill=cCyan, text=white, rounded corners=1.5pt, inner sep=3pt, font=\bfseries]{\thesection};}}
  {0.6em}{}
  [{\vspace{-0.3em}\color{cCyan}\titlerule[1.2pt]\vspace{0.3em}}]

\titleformat{\subsection}
  {\normalfont\large\bfseries\color{cTeal}}
  {\color{cTeal}\thesubsection}{0.6em}{}

\titlespacing*{\section}{0pt}{1.4em}{0.8em}
\titlespacing*{\subsection}{0pt}{1.1em}{0.5em}

% ---------- Cajas reutilizables ----------
\newtcolorbox{objetivobox}{
  colback=cBg, colframe=cCyan, boxrule=0.9pt, arc=2mm,
  left=8pt, right=8pt, top=6pt, bottom=6pt,
  title={\faBullseye\ \ Objetivo}, coltitle=white,
  colbacktitle=cTealDeep, fonttitle=\bfseries,
  before upper={\raggedright\color{cText}}
}

\newtcolorbox{incidenciabox}{
  colback=white, colframe=cRed, boxrule=1pt, arc=1mm,
  left=8pt, right=8pt, top=6pt, bottom=6pt,
  title={\faExclamationTriangle\ \ Hallazgo técnico}, coltitle=white,
  colbacktitle=cRed, fonttitle=\bfseries,
  before upper={\raggedright\color{cText}}
}

\newtcolorbox{evidenciabox}{
  colback=cBg, colframe=cMint, boxrule=0.9pt, arc=2mm,
  left=8pt, right=8pt, top=6pt, bottom=6pt,
  title={\faFileArchive\ \ Evidencia de ejecución}, coltitle=white,
  colbacktitle=cTealDeep, fonttitle=\bfseries,
  before upper={\raggedright\color{cText}}
}

% ---------- Figura sin marco ----------
\newcommand{\figuraMarcada}[3]{%
  \begin{figure}[H]
    \centering
    \includegraphics[width=#2\linewidth]{#1}
    \caption{#3}
  \end{figure}
}

\captionsetup{labelfont={bf,color=cTealDeep}, font=small, textfont=it}

% ============================================================
\begin{document}

% -------------------- PORTADA --------------------
\begin{titlepage}
\newgeometry{margin=0cm}
\begin{tikzpicture}[remember picture, overlay]
  \fill[cNavy] (current page.south west) rectangle (current page.north east);

  \fill[cPurple, opacity=0.28] (current page.north east) ++(-1,1) circle (3.2);
  \fill[cCyan, opacity=0.22] (current page.north east) ++(-2.6,0.4) circle (2.1);
  \fill[cTealDeep, opacity=0.55] (current page.south west) ++(1,-1) circle (4.0);
  \fill[cMint, opacity=0.16] (current page.south west) ++(2.4,-0.2) circle (2.3);

  % ---- Icono insignia ----
  \node[circle, fill=cTealDeep, draw=cCyan, line width=1.3pt, minimum size=2.5cm] (logo) at ($(current page.north)+(0,-2.9)$) {};
  \node at (logo) {\color{cMint}\fontsize{28}{28}\selectfont\faShoppingBasket};

  \node[anchor=north, text width=14cm, align=center] at ($(current page.north)+(0,-4.7)$) {
    {\color{cCyan}\bfseries\fontsize{12}{14}\selectfont \faDatabase\ \ MAESTRÍA EN CIENCIA DE DATOS}
  };

  \node[anchor=center, text width=15cm, align=center] at ($(current page.center)+(0,1.2)$) {
    {\color{white}\bfseries\fontsize{40}{44}\selectfont Práctica 5}\\[0.35em]
    {\color{cMint}\bfseries\fontsize{20}{24}\selectfont \faLink\ Reglas de Asociación con Apriori}
  };

  \draw[cCyan, line width=1.4pt] ($(current page.center)+(-3.6,-0.55)$) -- ++(7.2,0);

  \node[anchor=center, text width=13cm, align=center] at ($(current page.center)+(0,-1.3)$) {
    {\color{white!85!black}\fontsize{11}{14}\selectfont
     Itemsets frecuentes calculados a mano y reglas de asociación\\
     obtenidas con Apriori en Weka (minsup $\geq$ 0.5, minconf $\geq$ 0.6)}
  };

  \node[anchor=south, text width=14cm, align=center] at ($(current page.south)+(0,2.6)$) {
    \begin{tabular}{r l}
      \color{cCyan}\bfseries \faUser\ Elaborado por: & \color{white} Miguel Deza \\[3pt]
      \color{cCyan}\bfseries \faGraduationCap\ Programa: & \color{white} Maestría en Ciencia de Datos \\[3pt]
      \color{cCyan}\bfseries \faBook\ Curso: & \color{white} Introducción a Ciencia de Datos \\[3pt]
      \color{cCyan}\bfseries \faChalkboardTeacher\ Docente: & \color{white} Yessenia Deysi Yari Ramos \\[3pt]
      \color{cCyan}\bfseries \faIcon{calendar-alt}\ Fecha: & \color{white} 19 de septiembre de 2026 \\
    \end{tabular}
  };
\end{tikzpicture}
\end{titlepage}
\restoregeometry

% -------------------- 4. DESARROLLO --------------------
\section{\faCogs\ \ Desarrollo}

\subsection{\faSearch\ \ Dataset y exploración}
La tabla de partida fue la siguiente (``SI'' = el ítem está presente en la transacción):

\begin{table}[H]
\centering
\small
\begin{tabular}{ccccc}
\toprule
\textbf{Transacción} & \textbf{A} & \textbf{B} & \textbf{C} & \textbf{E} \\
\midrule
1 & SI & SI & SI & NO \\
2 & NO & SI & SI & SI \\
3 & SI & SI & SI & SI \\
4 & NO & SI & NO & SI \\
\bottomrule
\end{tabular}
\caption{Dataset original: 4 transacciones, 4 atributos nominales binarios.}
\end{table}

Los 4 atributos son \textbf{nominales} (\{SI,NO\}), sin valores faltantes, longitud fija de 4 instancias. Para cuantificar variabilidad codifiqué SI=1/NO=0 por atributo:

\begin{table}[H]
\centering
\small
\begin{tabular}{lrrrrr}
\toprule
\textbf{Atributo} & \textbf{Frec. SI} & \textbf{Frec. NO} & \textbf{Media} & \textbf{Var. muestral} & \textbf{Desv. Est.} \\
\midrule
A & 2 & 2 & 0.50 & 0.333 & 0.577 \\
B & 4 & 0 & 1.00 & 0.000 & 0.000 \\
C & 3 & 1 & 0.75 & 0.250 & 0.500 \\
E & 3 & 1 & 0.75 & 0.250 & 0.500 \\
\bottomrule
\end{tabular}
\caption{Distribución, media y desviación estándar por atributo (codificación SI=1/NO=0).}
\end{table}

\figuraMarcada{img/distribucion-atributos-weka.png}{0.95}{Distribución de valores por atributo en Weka (Preprocess $\to$ Visualize All): A balanceado 50/50, B constante, C y E con 75\% SI.}

\textbf{B es constante} (varianza 0): aparece en el 100\% de las transacciones, así que no discrimina nada y va a inflar artificialmente la confianza de cualquier regla donde sea consecuente. \textbf{A} es el atributo con más variabilidad (balance 50/50). En cuanto a relaciones entre pares, A y C coinciden en las dos transacciones donde A=SI (asociación fuerte real), mientras que C y E solo se solapan en 2 de las 3 transacciones donde aparecen (asociación más débil).

\subsection{\faLayerGroup\ \ Itemsets frecuentes (cálculo manual)}
Con soporte mínimo 0.5 (2/4 transacciones) y poda por la propiedad Apriori (si un subconjunto no es frecuente, el superconjunto tampoco), obtuve 11 itemsets frecuentes. \{A,E\} no es frecuente (soporte 0.25), así que todo itemset que lo contiene queda descartado:

\begin{table}[H]
\centering
\small
\begin{tabular}{clcc}
\toprule
\textbf{Nivel} & \textbf{Itemset} & \textbf{Soporte (conteo)} & \textbf{Soporte (\%)} \\
\midrule
$L_1$ & \{A\} & 2 & 0.50 \\
$L_1$ & \{B\} & 4 & 1.00 \\
$L_1$ & \{C\} & 3 & 0.75 \\
$L_1$ & \{E\} & 3 & 0.75 \\
$L_2$ & \{A,B\} & 2 & 0.50 \\
$L_2$ & \{A,C\} & 2 & 0.50 \\
$L_2$ & \{B,C\} & 3 & 0.75 \\
$L_2$ & \{B,E\} & 3 & 0.75 \\
$L_2$ & \{C,E\} & 2 & 0.50 \\
$L_3$ & \{A,B,C\} & 2 & 0.50 \\
$L_3$ & \{B,C,E\} & 2 & 0.50 \\
\bottomrule
\end{tabular}
\caption{Itemsets frecuentes por nivel (4 + 5 + 2 = 11 en total).}
\end{table}

\subsection{\faListOl\ \ Reglas de asociación}
De cada itemset frecuente de tamaño $\geq 2$ se generan todas las particiones antecedente$\to$consecuente, filtrando por confianza $\geq 0.6$. Resultan 19 reglas:

\begin{table}[H]
\centering
\small
\begin{minipage}{0.48\linewidth}
\centering
\begin{tabular}{lrr}
\toprule
\textbf{Regla} & \textbf{Sop.} & \textbf{Conf.} \\
\midrule
A $\to$ B    & 0.50 & 1.00 \\
A $\to$ C    & 0.50 & 1.00 \\
C $\to$ A    & 0.50 & 0.67 \\
B $\to$ C    & 0.75 & 0.75 \\
C $\to$ B    & 0.75 & 1.00 \\
B $\to$ E    & 0.75 & 0.75 \\
E $\to$ B    & 0.75 & 1.00 \\
C $\to$ E    & 0.50 & 0.67 \\
E $\to$ C    & 0.50 & 0.67 \\
A $\to$ BC   & 0.50 & 1.00 \\
\bottomrule
\end{tabular}
\end{minipage}%
\begin{minipage}{0.48\linewidth}
\centering
\begin{tabular}{lrr}
\toprule
\textbf{Regla} & \textbf{Sop.} & \textbf{Conf.} \\
\midrule
C $\to$ AB   & 0.50 & 0.67 \\
AB $\to$ C   & 0.50 & 1.00 \\
AC $\to$ B   & 0.50 & 1.00 \\
BC $\to$ A   & 0.50 & 0.67 \\
C $\to$ BE   & 0.50 & 0.67 \\
E $\to$ BC   & 0.50 & 0.67 \\
BC $\to$ E   & 0.50 & 0.67 \\
BE $\to$ C   & 0.50 & 0.67 \\
CE $\to$ B   & 0.50 & 1.00 \\
 & & \\
\bottomrule
\end{tabular}
\end{minipage}
\caption{Las 19 reglas de asociación con confianza $\geq 0.6$, generadas a partir de $L_2$ y $L_3$.}
\end{table}

\subsection{\faDesktop\ \ Ejecución en Weka}
Cargué el mismo dataset en Weka Explorer (pestaña \textbf{Preprocess}) y configuré Apriori en \textbf{Associate} con \texttt{lowerBoundMinSupport = 0.5}, \texttt{minMetric (confidence) = 0.6}, \texttt{numRules = 100} (para no cortar antes de llegar al soporte mínimo) y \texttt{outputItemSets = True}.

\figuraMarcada{img/evidencia-dataset-cargado.png}{0.92}{Dataset cargado en Weka Explorer: 4 atributos nominales (A, B, C, E), sin valores nulos.}

\figuraMarcada{img/resultado-apriori-itemsets.png}{0.92}{Salida de Apriori: itemsets frecuentes por nivel -- $L_1{=}5$, $L_2{=}7$, $L_3{=}3$.}

\figuraMarcada{img/resultado-apriori-reglas.png}{0.85}{Reglas de asociación generadas por Apriori en Weka (soporte, confianza y lift).}

\end{document}
